\section{Sufism}

Sufism is the mystical and contemplative tradition within Islam. It holds the
strict oneness of God shared by all of Islam, but develops it inward, treating
the cosmos as the self-disclosure of the one Real and the soul's task as the
return to that source. Its most systematic statement comes from Ibn Arabi, and
its practice centers on the refinement of the lower self and the stages of union.

\subsection{The Creation Model}

In the beginning there is one Real. It is utterly one, alone, and beyond all
likeness. In itself it is a hidden treasure, complete and unknown. It wishes to
be known, and out of that wish the world arises. The one breathes its own names
outward, and each name takes form as a creature. Every thing that exists is the
trace of a name, the outward showing of one face of the hidden one. So the
cosmos is not a thing standing apart from the source. It is a mirror. In it the
one sees its own names reflected back, and by seeing them it knows itself.

The breath that gives existence is mercy. To exist at all is to be touched by it.
This breathing is not a single past act. The world is re-breathed at every
instant, made new moment by moment. Among all creatures the human being is the
clearest part of the mirror, the one being that can hold all the names at once.
In the perfected human the reflection is complete, and the one beholds itself
fully. The path of the soul is to lose its separate self in the one and then to
abide in the one, so that the wave returns to the sea it never truly left.

\subsection{Oneness and the Hidden Essence}

The defining move is the assertion of absolute oneness (\textbf{tawhid}). Nothing
shares in the divine. The Sufis prefer to call the absolute the Real
(\textbf{al-Haqq}). The \textbf{Essence} (dhat) is God in himself, the Hidden (al-Ghayb),
never reached even by the prophets. What creation can know are the \textbf{99
names} (al-asma al-husna), the faces through which the Essence turns toward the
world. The names split into two families. The \textbf{jamal} names disclose
gentleness and nearness, such as the Merciful (al-Rahman) and the Loving
(al-Wadud). The \textbf{jalal} names disclose majesty and rigor, such as the
Compeller (al-Jabbar) and the Subduer (al-Qahhar). The name \textbf{Allah} gathers
all of them.

\subsection{The Oneness of Being}

\textbf{Ibn Arabi} teaches the oneness of being (\textbf{wahdat al-wujud}). All
that exists is the one Being disclosing itself. Existence (wujud) is a single
reality, and creatures are its determinations. The mechanism is the
\textbf{breath of the Merciful} (\textbf{nafas al-Rahman}). As speech is breath shaped by
the mouth, the cosmos is the breath of God shaped by the names. This breath is
continuous, so creation is renewed at every instant.

Before things exist outwardly they exist as archetypes in the divine knowledge.
These are the \textbf{fixed entities} (\textbf{a'yan thabita}). Each thing
has its own fixed entity, known to God from eternity. They have no being of
their own. They are the receptacles into which existence is poured. A thing's
nature is set by its archetype, not imposed from outside. God gives only
existence, and the thing brings its character from its own fixed entity.

The unifying idea is self-disclosure (\textbf{tajalli}). The hidden treasure
wills to be known, discloses itself through the names, and the cosmos is the
\textbf{mirror} in which the names appear. The \textbf{perfect man}
(\textbf{al-insan al-kamil}) is the most polished part of that mirror, the point at
which the reflection is complete and God sees himself fully.

\subsection{The Outer Frame of the Five Pillars}

The inner path of the Sufi sits within the outer practice shared by all of
Islam, the \textbf{five pillars}. They are the public frame on which the
contemplative life rests.

\begin{longtable}{L{0.06\linewidth} L{0.24\linewidth} L{0.54\linewidth}}
\caption{The five pillars of Islam.}\\
\toprule
\# & Pillar & Meaning \\
\midrule
1 & shahada & the testimony that there is no god but God and Muhammad is his messenger \\
2 & salat & the five daily prayers \\
3 & zakat & the giving of alms \\
4 & sawm & the fast of Ramadan \\
5 & hajj & the pilgrimage to Mecca \\
\bottomrule
\end{longtable}

The Sufis keep the pillars and read each one inwardly, as an outward act that
shapes the heart. The pillars are the law (shariah), and the inner path
(tariqa) travels through them toward the Real (haqiqa).

\subsection{The Seven Valleys of the Path}

The poet \textbf{Attar}, in the Conference of the Birds, maps the journey to God
as a passage through \textbf{seven valleys}. The birds cross them in order to
reach their king.

\begin{longtable}{L{0.06\linewidth} L{0.28\linewidth} L{0.50\linewidth}}
\caption{The seven valleys of Attar.}\\
\toprule
\# & Valley & Meaning \\
\midrule
1 & seeking (talab) & the search begins, old certainties are left behind \\
2 & love (ishq) & burning desire for the one, beyond reason \\
3 & knowledge (marifa) & direct knowing, each soul by its own way \\
4 & detachment (istighna) & independence, the dropping of worldly claims \\
5 & unity (tawhid) & all is seen as one, the many fall away \\
6 & wonderment (hayrat) & bewilderment, knowing and not knowing at once \\
7 & poverty and annihilation (faqr and fana) & the self is spent and passes away into the one \\
\bottomrule
\end{longtable}

\subsection{The Soul and Its Refinement}

The lower self (the \textbf{nafs}) is refined as it ascends. The fullest scheme
gives seven stages, each named from a phrase in the Quran.

\begin{longtable}{L{0.06\linewidth} L{0.34\linewidth} L{0.46\linewidth}}
\caption{The seven stages of the nafs, the lower self.}\\
\toprule
\# & Stage & Meaning \\
\midrule
1 & al-nafs al-ammara & The soul that commands to evil \\
2 & al-nafs al-lawwama & The self-reproaching soul \\
3 & al-nafs al-mulhima & The inspired soul \\
4 & al-nafs al-mutmainna & The soul at peace \\
5 & al-nafs al-radiyya & The well-pleased soul \\
6 & al-nafs al-mardiyya & The soul pleasing to God \\
7 & al-nafs al-safiyya & The pure or perfected soul \\
\bottomrule
\end{longtable}

Sufi psychology also names \textbf{subtle centers} (lataif), inner organs of
perception. The common scheme counts the heart (\textbf{qalb}), the spirit
(\textbf{ruh}), the secret (\textbf{sirr}), the hidden (\textbf{khafi}), and the
most hidden (\textbf{akhfa}), with the body holding them all. The number and
the order vary by order.

The final movements of the path are \textbf{fana} and \textbf{baqa}.
\textbf{Fana} is annihilation, the passing-away of the self in God. The drop
vanishes into the sea, and the seeker no longer sees a self apart from the one.
\textbf{Baqa} follows. The self does not stay annihilated. It abides in God,
now acting through the divine rather than the ego. The order is fana then baqa,
to die before dying and then to live the divine life.

\subsection{The Worlds}

The cosmos descends in planes, from pure divinity to the senses.

\begin{longtable}{L{0.20\linewidth} L{0.62\linewidth}}
\caption{The descent of the worlds from divinity to the senses.}\\
\toprule
World & Nature \\
\midrule
lahut & Divinity, the realm of the Essence and the names \\
jabarut & Power, the realm of pure spirits and angels \\
malakut & The unseen, the realm of souls and images \\
mulk & The visible, the realm of bodies and the senses \\
\bottomrule
\end{longtable}

The threefold core is mulk, malakut, and jabarut. Fourfold schemes add lahut on
top, and fivefold schemes add hahut, the bare Essence, above all. The exact
naming varies by text.

\subsection{Comparison to the Base Models}

The structure matches the base models closely. A hidden source beyond name
discloses itself through ranked degrees of being, and the soul climbs back
through stages toward union. The hidden treasure that wished to be known plays
the role of the overflowing first principle. The fixed entities parallel the
archetypes that pattern the lower worlds. Fana and baqa are the
go-out-and-return read at the level of the self, the wave dissolving into the
sea and then abiding in it. The breath goes out as the names and returns as the
soul that knows the one.
